\documentclass{article}
\usepackage[utf8]{inputenc}    
\usepackage[T1]{fontenc}       
\usepackage{lmodern}           
\usepackage{hyperref}
\hypersetup{
    colorlinks=true,       % false: boxed links; true: colored links
    linkcolor=blue,        % color of internal links (sections, etc.)
    urlcolor=cyan,         % color of external links
    linktoc=all            % 'all' makes both text and page number clickable
}
\usepackage{bbm}
\usepackage{amsmath}   
\usepackage{amssymb}   
\usepackage{geometry}  
\usepackage{enumerate} 
\usepackage{xcolor}  
\usepackage{amsthm}
\usepackage{pdfpages}
\usepackage{mathrsfs}
\usepackage{stmaryrd}
\newtheorem{theorem}{Theorem}[section]
% \newtheorem{lemma}[theorem]{Lemma}
\newtheorem{corollary}[theorem]{Corollary}
% \newtheorem{definition}[theorem]{Definition}
\newtheorem{example}[theorem]{Example}
\usepackage{listings}  
\usepackage{tikz-cd}
\usetikzlibrary{positioning}
\usepackage{forest}     
\usepackage{xcolor}
\usepackage{tcolorbox}
\tcbuselibrary{theorems}
\usepackage{algpseudocode}
\usepackage[ruled, vlined, algoruled]{algorithm2e} % Options for styling
\usepackage{dsfont}
% remove/comment out the old line:
% \newtheorem{definition}[theorem]{Definition}

\usepackage{xcolor}
\usepackage{tcolorbox}
\tcbuselibrary{breakable}
\newtheorem{definition}[theorem]{Definition}
\let\olddefinition\definition
\let\endolddefinition\enddefinition
\renewenvironment{definition}
  {\begin{tcolorbox}[
      colback=red!10,
      colframe=red!50!black,
      boxrule=0.6pt,
      arc=2pt,
      left=6pt,right=6pt,top=4pt,bottom=4pt,
      breakable
    ]\olddefinition}
  {\endolddefinition\end{tcolorbox}}

\newtheorem{lemma}[theorem]{Lemma}
\let\oldlemma\lemma
\let\endoldlemma\endlemma
\renewenvironment{lemma}
  {\begin{tcolorbox}[
      colback=green!8,
      colframe=green!50!black,
      boxrule=0.6pt,
      arc=2pt,
      left=6pt,right=6pt,top=4pt,bottom=4pt
    ]\oldlemma}
  {\endoldlemma\end{tcolorbox}}

  \let\oldtheorem\theorem
\let\endoldtheorem\endtheorem
\renewenvironment{theorem}
  {\begin{tcolorbox}[
      colback=yellow!6,
      colframe=yellow!60!black,
      boxrule=0.6pt,
      arc=2pt,
      left=6pt,right=6pt,top=4pt,bottom=4pt,
      parbox=false
    ]\oldtheorem}
  {\endoldtheorem\end{tcolorbox}}
\usepackage{bussproofs}
\usetikzlibrary{arrows.meta}
\usetikzlibrary{arrows.meta,decorations.pathreplacing,calc}
\lstset{frame=tb,
  language=C,
  aboveskip=3mm,
  belowskip=3mm,
  showstringspaces=false,   
  columns=flexible,
  basicstyle={\small\ttfamily},
  numbers=none,
  numberstyle=\tiny\color{gray},
  keywordstyle=\color{blue},
  commentstyle=\color{brown},
  stringstyle=\color{orange},
  breaklines=true,
  breakatwhitespace=true,
  tabsize=3
}
\geometry{top=0.5in, bottom=0.5in, left=1in, right=1in}
% ctrl shift D to toggle night 
\begin{document}

\title{HW4}
\author{Wang Xiyu A0282500R}
\date{}
\maketitle
\section*{Skolem Normal Form and Skolemization}

\begin{definition}[Skolem Normal Form]
A formula is in Skolem Normal Form if it is a sentence in prenex normal form with no existential quantifiers.
\end{definition}

\begin{definition}[Skolemization]
Let $\varphi$ be a FOL formula. Its Skolem Normal Form $\mathrm{sk}(\varphi)$ is obtained as follows:
\begin{enumerate}
    \item Convert $\varphi$ to an equivalent prenex normal form $\psi$ over the same signature.
    \item Existentially quantify all free variables to obtain a sentence $\psi_0$.
    \item Replace each existential variable $x$ by a Skolem term:
    \begin{itemize}
        \item If $\mathrm{dep}(x) = \emptyset$, replace $x$ by a fresh constant $c_x$.
        \item Otherwise, replace $x$ by a fresh function $f_x(\bar y)$ where $\bar y = \mathrm{dep}(x)$.
    \end{itemize}
    Remove all existential quantifiers.
\end{enumerate}
\end{definition}

\paragraph{General Form}
\[
\psi_0 \equiv \exists x_{1,1}\cdots \exists x_{1,k_1} \forall y_{1,1}\cdots \forall y_{1,l_1} \cdots
\exists x_{m,1}\cdots \exists x_{m,k_m} \forall y_{m,1}\cdots \forall y_{m,l_m} \;\psi_b
\]

\[
\mathrm{sk}(\varphi) =
\forall y_{1,1}\cdots \forall y_{1,l_1}\cdots \forall y_{m,1}\cdots \forall y_{m,l_m}\;
\psi_b[x_{i,j} \mapsto \mathrm{sTrm}(x_{i,j})]
\]

\paragraph{Properties}
\begin{itemize}
    \item $\mathrm{sk}(\varphi)$ is a sentence
    \item $\mathrm{sk}(\varphi)$ has no existential quantifiers
    \item $\mathrm{sk}(\varphi)$ is over an extended signature
\end{itemize}


\section*{Question 1}
Show that for any FOL formula $\varphi$,
\[
\varphi \text{ is satisfiable } \iff \mathrm{sk}(\varphi) \text{ is satisfiable}.
\]

\begin{proof}
    Consider skolemization as a recursive algorithm described as follows:\\
\begin{algorithm}[H]
    \caption{Skolemization on Prenex Formula}
    \KwIn{
    Signature $\Sigma = (C, F, R)$; \\
    Formula $\phi = (\pi, \psi)$ where $\pi$ is a quantifier prefix and $\psi$ is quantifier-free\;
    List $u$ of universally quantified variables in scope (initially empty).
    }
    \KwOut{
    Skolemized formula $(\pi', \psi')$ equisatisfiable with $\phi$.
    }
    
    \If{$\pi$ is empty}{
        \Return$(\emptyset, \psi)$
    }
    
    $(q, x) \gets \text{head}(\pi)$\;
    $\pi_{\text{rest}} \gets \text{tail}(\pi)$
    
    \If{$q = \exists$}{
        \eIf{$u$ is empty}{
            introduce fresh constant $c_x$\;
            $\Sigma.C \gets \Sigma.C \uplus \{c_x\}$\;
            \Return{Skol$(\Sigma, (\pi_{\text{rest}}, \psi[x \mapsto c_x]), u)$}
        }{
            introduce fresh function $f_x$ of arity $|u|$\;
            $\Sigma.F \gets \Sigma.F \uplus \{f_x\}$\;
            let $t = f_x(u_1, \dots, u_k)$ where $u = [u_1, \dots, u_k]$\;
            \Return{Skol$(\Sigma, (\pi_{\text{rest}}, \psi[x \mapsto t]), u)$}
        }
    }
    \Else{ \tcp{$q = \forall$}
        $u' \gets u \cdot x$ \tcp{append $x$ to $u$}
        $(\pi', \psi') \gets$ Skol$(\Sigma, (\pi_{\text{rest}}, \psi), u')$\;
        \Return{$((\forall x)\cdot \pi', \psi')$}
    }
\end{algorithm}
We claim that the invariant is the satisfiability of input and output is the same.
Let $\text{Skol}(\Sigma, (\pi, \psi), u)$ be the recursive algorithm defined. 
We prove by induction on the length of the quantifier prefix $n = |\pi|$ that the algorithm preserves equisatisfiability.
\subsubsection*{Inductive hypotshis}
For any prefix $\pi$ of length $n$, any signature $\Sigma$, and any list of universal variables $u = [u_1, \dots, u_k]$, 
a formula $\phi = \forall u (\pi \psi)$ is satisfiable in a $\Sigma$-structure $\mathfrak{A}$ 
if and only if there exists a larger $\mathfrak{A}'$ than $\mathfrak{A}$ to $\Sigma'$ (containing new symbols) 
such that $\mathfrak{A}'$ satisfies the output of $\text{Skol}(\Sigma, (\pi, \psi), u)$.
\[\mathfrak{A} \models \phi = \forall u (\pi \psi) \iff \mathfrak{A}' \models {Skol}(\Sigma, (\pi, \psi), u)\]

\subsection*{Case 1: Base Case ($n = 0$)}
If $\pi$ is empty, the algorithm returns $(\emptyset, \psi)$, trivially equisatisfiable

\subsection*{Case 2: Universal Case ($q = \forall$)}
Let $\pi = \forall x \cdot \pi_{\text{rest}}$. The algorithm sets $u' = u \cdot x$ and recurses\\
We will show that:
$$\mathfrak{A} \models \forall u \forall x (\pi_{\text{rest}} \psi) \iff \mathfrak{A}' \models \forall u \forall x ({Skol}(\pi_{\text{rest}}, \psi, u'))$$

\begin{itemize}
    \item \textbf{ ($\Rightarrow$):} Suppose $\mathfrak{A} \models \forall u \forall x (\pi_{\text{rest}} \psi)\implies \mathfrak{A} \models \forall (u \cup \{x\}) (\pi_{\text{rest}} \psi)$. 
    There exists an expansion $\mathfrak{A}'$ such that $\mathfrak{A}' \models \forall (u \cup \{x\}) {Skol}(\pi_{\text{rest}}, \psi, u')$. rearrange the quantifiers gives $\mathfrak{A}' \models \forall u \forall x {Skol}(\dots)$.
    \item \textbf{ ($\Leftarrow$):} Suppose $\mathfrak{A}' \models {Skol}(\pi_{\text{rest}}, \psi, u')$. Since $\mathfrak{A}'$ is an expansion of $\mathfrak{A}$, and the universal quantifier $\forall x$ 
    does not introduce new symbols in this step, same logic 
\end{itemize}

\subsection*{Case 3: Existential Case ($q = \exists$)}

\end{proof}


\section*{Theorem 1 (Restricted Downward Löwenheim-Skolem)}
Let $\Sigma$ be a countable signature and $\Gamma$ a set of universally quantified FOL sentences without equality.  
If $\Gamma$ is satisfiable, then there exists a countable structure $\mathfrak{A}$ such that
\[
\mathfrak{A} \models \Gamma.
\]

\section*{Theorem 2 (Full Downward Löwenheim-Skolem)}
Let $\Sigma$ be a countable signature, $V$ a countable set of variables, and $\Gamma$ a set of FOL formulae.  
If $\Gamma$ is satisfiable, then there exists a countable structure $\mathfrak{A}$ such that
\[
\mathfrak{A} \models \Gamma.
\]

\section*{Question 2}
Prove Theorem 2.
\begin{proof}
    Let $\Gamma$ be a set of FOL formulae over a countable signature $\Sigma$.
    \begin{enumerate}
        \item For each $\phi \in \Gamma$ with free variables $x_1, \dots, x_n$, let $\phi' = \exists x_1 \dots \exists x_n \phi$. $\Gamma$ is sat $\iff \Gamma' = \{ \phi' \mid \phi \in \Gamma \}$ is satisfiable
        \item Convert $\Gamma'$ to $sk(\Gamma)$ by skolmization. 
        $sk(\Gamma)$ is a set of universal sentences over an extended countable signature $\Sigma'$, By Question 1, $sk(\Gamma)$ is sat $\iff \Gamma'$ is sat
        \item Since $sk(\Gamma)$ consists of universal sentences, by Theorem 1, there exists a countable structure $\mathfrak{A} = (U, I)$ such that $\mathfrak{A} \models sk(\Gamma)$
        \item If $\Sigma$ contains equality, treat $=$ as a binary relation $EQ$. By the congruence proof from hw4, define $I(EQ)$ to make a congruence relation on $\mathfrak{A}$. 
        \item construct a new quotient structure $\mathfrak{A}_{EQ}$ with universe $U_{EQ} = \{ [u]_{EQ} \mid u \in U \}$ where equal elemnts collapse to th same equivalent class
        \item $|U|$ is countable $\implies |U/EQ| \le |U|$ is countable. By congruence, $\mathfrak{A}/EQ \models sk(\Gamma) \implies \mathfrak{A}/EQ \models \Gamma$.
    \end{enumerate}
    \end{proof}

\section*{Theorem 3 (Compactness)}
Let $\Sigma$ be a countable signature and $\Gamma$ a set of FOL formulae. Then:
\[
\Gamma \text{ is satisfiable } \iff \Gamma \text{ is finitely satisfiable}.
\]

\section*{Question 3}
Prove Theorem 3.

\section*{Definability}

\begin{definition}
Let $\Sigma$ be a signature and $C$ a class of $\Sigma$-structures.  
$C$ is definable in FOL if there exists a sentence $\varphi$ such that
\[
\llbracket \varphi \rrbracket = \{ \mathfrak{A} \mid \mathfrak{A} \models \varphi \} = C.
\]
\end{definition}

\section*{Question 4 (Reachability)}
Prove that reachability is not definable in FOL, i.e., there is no sentence $\varphi_{\text{reach}}$ such that
\[
\{ G(\mathfrak{A}) \mid \mathfrak{A} \models \varphi_{\text{reach}} \} = G_{\text{reach}}.
\]

\begin{proof}[Proof of Question 4]
    To prove that reachability is not definable in First-Order Logic (FOL), we proceed by contradiction using the Compactness Theorem
    
    \begin{enumerate}
        \item
        Let $\Sigma_{graphs}$ be the signature $( \{s, t\}, \emptyset, \{E\} )$
        For each $k \in \mathbb{N}$, define a FOL formula $\psi_k(s, t)$ asserting there exists a path of length at most $k$ from $s$ to $t$:
        \[ \psi_k(s, t) \equiv \exists v_0 \dots \exists v_k \left( (v_0 = s) \wedge (v_k = t) \wedge \bigwedge_{0 \le i < k} E(v_i, v_{i+1}) \right) \]
    
        \item 
        Suppose there exists a sentence $\phi_{reach}$ such that $\mathfrak{A} \models \phi_{reach}$ iff $t$ is reachable from $s$ in $G(\mathfrak{A})$
        Consider the set of sentences:
        \[ \Gamma = \{ \phi_{reach} \} \cup \{ \neg \psi_k(s, t) \mid k \in \mathbb{N} \} \]
    
        \item 
        Let $\Gamma_0 \subseteq \Gamma$ be any finite subset. There exists a maximum $k_{max}$ such that $\neg \psi_{k_{max}}$ is in $\Gamma_0$.
        Construct a graph $G$ where $t$ is reachable from $s$ via a path of length $k_{max} + 1$. 
        In this structure, $\phi_{reach}$ is true, and all $\neg \psi_k$ for $k \le k_{max}$ are true. 
        Thus, $\Gamma_0$ is finitely satisfiable.
    
        \item
        By the Compactness Theorem, if every finite subset of $\Gamma$ is satisfiable, then $\Gamma$ is satisfiable
        However, a model $\mathfrak{A} \models \Gamma$ would imply:
        \begin{itemize}
            \item $t$ is reachable from $s$ (since $\mathfrak{A} \models \phi_{reach}$)
            \item There is no finite path from $s$ to $t$ (since $\mathfrak{A} \models \neg \psi_k$ for all $k$)
        \end{itemize}
        contradiction, as reachability requires a finite sequence of vertices
        Therefore, no such sentence $\phi_{reach}$ exists.
    \end{enumerate}
    \end{proof}

\section*{Question 5 ($k$-Colorability)}
Define a FOL sentence $\varphi_{k\text{-colorable}}$ such that
\[
\{ G(\mathfrak{A}) \mid \mathfrak{A} \models \varphi_{k\text{-colorable}} \} = G_{k\text{-colorable}}.
\]

\begin{proof}
    To define $k$-colorability over the signature $\Sigma_{colored-graphs} = (\emptyset, \emptyset, \{E, R_1, \dots, R_k\})$, we construct a sentence $\phi_{k-colorable}$ that enforces the following constraints:
    
    \begin{enumerate}
        \item {Total Assignment:} 
        \[ \forall x  \bigvee_{i=1}^k R_i(x)  \]
        
        \item {Mutual Exclusivity:}
        \[ \forall x  \bigwedge_{1 \le i < j \le k} \neg (R_i(x) \wedge R_j(x))  \]
        
        \item {Edge Consistency:} 
        \[ \forall x \forall y E(x, y) \implies \bigwedge_{i=1}^k \neg (R_i(x) \wedge R_i(y))  \]
    \end{enumerate}
    
    \[ \phi_{k-colorable} \equiv \forall x \left( \bigvee_{i=1}^k R_i(x) \right) \wedge \forall x \left( \bigwedge_{i < j} \neg (R_i(x) \wedge R_j(x)) \right) \wedge \forall x \forall y \left( E(x, y) \implies \bigwedge_{i=1}^k \neg (R_i(x) \wedge R_i(y)) \right) \]
    \end{proof}

\begin{definition}
    Let $\Sigma = (C, F, R)$ and $\mathfrak{A}_1, \mathfrak{A}_2$ be $\Sigma$-structures.  
A function $h : U_1 \to U_2$ is a homomorphism if:
\begin{enumerate}
    \item $h(I_1(c)) = I_2(c)$ for all $c \in C$
    \item $h(I_1(f)(\bar u)) = I_2(f)(h(\bar u))$ for all $f \in F$
    \item $(\bar u \in I_1(R)) \iff (h(\bar u) \in I_2(R))$ for all $R \in R$
\end{enumerate}
\end{definition}


\section*{Question 6}
\begin{enumerate}
    \item Prove $h(n) = n \bmod 2$ is a homomorphism from $\mathfrak{A}_\mathbb{N}$ to $\mathfrak{A}_\mathbb{B}$.
    \item Find $\varphi$ such that $\mathfrak{A}_\mathbb{N} \models \varphi$ but $\mathfrak{A}_\mathbb{B} \not\models \varphi$.
\end{enumerate}
6.1
\begin{proof}
    To prove $h: \mathbb{N} \to \mathbb{B}$ is a homomorphism from $\mathfrak{A}_{\mathbb{N}}$ to $\mathfrak{A}_{\mathbb{B}}$:
    \begin{enumerate}
        \item {Constants:} 
        \begin{itemize}
            \item $h(I_{\mathbb{N}}(\mathbb{O})) = h(0) = 0 = I_{\mathbb{B}}(\mathbb{O})$.
            \item $h(I_{\mathbb{N}}(\mathds{1})) = h(1) = 1 = I_{\mathbb{B}}(\mathds{1})$.
        \end{itemize}
        \item {Functions ($+$):} We must show $h(m + n) = I_{\mathbb{B}}(+)(h(m), h(n))$ for all $m, n \in \mathbb{N}$.
        \begin{itemize}
            \item Case 1: $m, n$ both even. $m+n$ is even. $h(m+n)=0$. $I_{\mathbb{B}}(+)(0, 0)=0$.
            \item Case 2: $m, n$ both odd. $m+n$ is even. $h(m+n)=0$. $I_{\mathbb{B}}(+)(1, 1)=0$.
            \item Case 3: One even, one odd. $m+n$ is odd. $h(m+n)=1$. $I_{\mathbb{B}}(+)(0, 1)=1$.
        \end{itemize}
    \end{enumerate}
    \end{proof}
    6.2
    \begin{proof}
    \[ \phi \equiv \neg (1 + 1 = 0) \]
    \begin{itemize}
        \item $1 + 1 = 2 \neq 0 \implies \mathfrak{A}_{\mathbb{N}} \models \phi$.
        \item $1 + 1 = 0 \implies \mathfrak{A}_{\mathbb{B}} \models \neg \phi$.
    \end{itemize}
    \end{proof}

\begin{definition}
    A homomorphism $h : U_1 \to U_2$ is an isomorphism if it is bijective.

\[
\mathfrak{A}_1 \cong \mathfrak{A}_2 \iff \exists h \text{ isomorphism}
\]
\end{definition}

\section*{Question 7}
Prove that if $\mathfrak{A}_1 \cong \mathfrak{A}_2$, then $\mathfrak{A}_2 \cong \mathfrak{A}_1$.
\begin{proof}
    Let $h: U_1 \to U_2$ be an isomorphism from $\mathfrak{A}_1$ to $\mathfrak{A}_2$. By definition, $h$ is a bijective homomorphism. Since $h$ is a bijection, there exists an inverse function $h^{-1}: U_2 \to U_1$ which is also a bijection. We show $h^{-1}$ is a homomorphism:
    
    \begin{enumerate}
        \item $\forall c \in {C}$, we know $h(I_1(c)) = I_2(c)$. Applying $h^{-1}$ to both sides, we get $I_1(c) = h^{-1}(I_2(c))$.
        
        \item $\forall f \in {F}$ of arity $k$ and $v_1, \dots, v_k \in U_2$, let $u_i = h^{-1}(v_i)$. 
        \[ h(I_1(f)(u_1, \dots, u_k)) = I_2(f)(h(u_1), \dots, h(u_k)) = I_2(f)(v_1, \dots, v_k) \]
        Applying $h^{-1}$ on both sides:
        \[ I_1(f)(h^{-1}(v_1), \dots, h^{-1}(v_k)) = h^{-1}(I_2(f)(v_1, \dots, v_k)) \]
        
        \item $\forall R \in \mathcal{R}$ of arity $k$:
        \[ (v_1, \dots, v_k) \in I_2(R) \iff (h(h^{-1}(v_1)), \dots, h(h^{-1}(v_k))) \in I_2(R) \]
        \[ \iff (h^{-1}(v_1), \dots, h^{-1}(v_k)) \in I_1(R) \]
    \end{enumerate}
    Since $h^{-1}$ is a bijective homomorphism, it is an isomorphism from $\mathfrak{A}_2$ to $\mathfrak{A}_1$.
    \end{proof}

\section*{Question 8}
\begin{enumerate}
    \item Give examples of isomorphic structures.
    \item Give examples of non-isomorphic structures with same finite cardinality.
\end{enumerate}

8.1 
\begin{proof}
Natural number to even natural number, shown in class
\end{proof}
8.2
\begin{proof}
    $\mathfrak{A}_1$ and  $\mathfrak{A}_2$ from question 6, just change $U_{\mathfrak{A}_1}$ from $\mathbb{N}$ to $\{0,1\}$, redefine $+$ such that $1+1 = 1$
\end{proof}

\section*{Theorem 4 (Isomorphism Theorem)}
Let $h : \mathfrak{A}_1 \to \mathfrak{A}_2$ be an isomorphism. Then for every formula $\varphi$ and assignment $\alpha$,
\[
\mathfrak{A}_1, \alpha \models \varphi \iff \mathfrak{A}_2, (h \circ \alpha) \models \varphi.
\]


  \end{document}